% 彗星（综述）
% license CCBYSA3
% type Wiki

本文根据 CC-BY-SA 协议转载翻译自维基百科\href{https://en.wikipedia.org/wiki/Comet}{相关文章}。

\begin{figure}[ht]
\centering
\includegraphics[width=8cm]{./figures/45ca94c34375617b.png}
\caption{} \label{fig_Comet_1}
\end{figure}
彗星是一类由冰组成的、体积较小的太阳系天体或星际天体。当其经过接近太阳的区域时，会因受热而开始释放气体，这一过程称为释气作用（outgassing）。这一过程会在彗核周围产生一个延展的、未被引力束缚的大气层，即彗发（coma），并有时形成由彗发中吹离的气体和尘埃构成的彗尾。这些现象源于太阳辐射以及向外流动的太阳风等离子体对彗核的作用。彗核的大小从数百米到数十公里不等，由松散聚集的冰、尘埃和小型岩石颗粒组成。彗发的尺度可达地球直径的 15 倍，而彗尾的长度甚至可能超过一个天文单位。如果距离足够近且亮度足够高，人类在无望远镜的情况下即可从地球上观测到彗星，其在天空中的角距可达 30°（相当于 60 个月亮）。自古以来，彗星已被众多文化和宗教观察并记录。

彗星通常具有高离心率的椭圆轨道，其轨道周期范围极广，从数年到可能长达数百万年。短周期彗星来源于海王星轨道之外的柯伊伯带或相关的散射盘。长周期彗星则被认为起源于奥尔特云，一个由冰质天体组成的球形云层，其范围从柯伊伯带外侧延伸至通往最近恒星方向一半的距离\(^\text{[2]}\)。长周期彗星通常因经过恒星的引力扰动或银河潮汐效应而被触发向太阳方向运动。双曲轨道彗星可能只会一次穿越内太阳系，随后被抛入星际空间。彗星的出现称为一次显现（apparition）。

多次接近太阳的灭绝彗星会失去几乎全部挥发性冰与尘埃，并可能逐渐呈现出类似小型小行星的外观\(^\text{[3]}\)。一般认为，小行星与彗星的起源不同：小行星形成于木星轨道以内，而非太阳系外部\(^\text{[4][5]}\)。然而，主带彗星（main-belt comets）以及活跃的半人马族小天体的发现，使小行星与彗星之间的界限变得模糊。21 世纪初，一些具备长周期彗星轨道却呈现内太阳系小行星特征的小天体被称为曼岛彗星（Manx comets）。它们仍被归类为彗星，例如 C/2014 S3（PANSTARRS）\(^\text{[6]}\)。在 2013 年至 2017 年之间，共发现了 27 颗曼岛彗星\(^\text{[7]}\)。

截至 2021 年 11 月，共已知 4,584 颗彗星\(^\text{[8]}\)。然而，这仅占可能存在的彗星总量的一小部分，因为太阳系外部（奥尔特云）中类彗星天体的储库估计约为一万亿个\(^\text{[9][10]}\)。平均每年约有一颗彗星可被肉眼看见，尽管其中许多较为暗淡、外观并不特别醒目\(^\text{[11]}\)。特别明亮的彗星被称为“伟大彗星”（great comets）。人类已利用无人探测器访问过彗星，例如美国 NASA 的“深度撞击号”（Deep Impact），其通过撞击坦普尔 1 号彗星形成陨坑以研究其内部结构；又如欧洲航天局的“罗塞塔号”（Rosetta），它成为首个在彗星上着陆的机器人探测器\(^\text{[12]}\)。
\subsection{词源}
\begin{figure}[ht]
\centering
\includegraphics[width=8cm]{./figures/6a3173d4cf8911ea.png}
\caption{公元 729 年的《盎格鲁－撒克逊编年史》以及尊敬的比德（the Venerable Bede）的著作中均提到了彗星。} \label{fig_Comet_2}
\end{figure}
单词“comet（彗星）”源自古英语 cometa，其来源为拉丁语 comēta 或 comētēs，而这些拉丁词又是对希腊语 κομήτης 的罗马化形式，意为“留着长发的（事物）”。《牛津英语词典》指出，术语（ἀστὴρ） κομήτης 在希腊语中早已表示“长发之星，即彗星”。κομήτης 源自动词 κομᾶν（koman），意为“留长发”，而该动词又派生自名词 κόμη（komē），意为“头发”，并被用来表示“彗尾”\(^\text{[15][16]}\)。

彗星的天文学符号（在 Unicode 中表示）为 U+2604 ☄ COMET，由一个小圆盘及其后三条类似发丝的延伸部分组成\(^\text{[17]}\)。
\subsection{物理特征}
\begin{figure}[ht]
\centering
\includegraphics[width=8cm]{./figures/37fee7dd52548b31.png}
\caption{彗星的结构} \label{fig_Comet_3}
\end{figure}
\subsubsection{彗核}
\begin{figure}[ht]
\centering
\includegraphics[width=8cm]{./figures/43510b130873771c.png}
\caption{在航天器飞越期间拍摄的 103P/Hartley 彗核图像。该彗核长度约为 2 千米。} \label{fig_Comet_4}
\end{figure}
彗星的固体核心结构被称为彗核。彗核由岩石、尘埃、水冰以及冻结的二氧化碳、一氧化碳、甲烷和氨等物质的混合体构成\(^\text{[18]}\)。因此，根据弗雷德·惠普尔（Fred Whipple）的模型，它们常被形象地称为“肮脏的雪球”（dirty snowballs）\(^\text{[19]}\)。尘埃含量更高的彗星则被称为“冰质尘球”（icy dirtballs）\(^\text{[20]}\)。“冰质尘球”这一术语源于 2005 年 7 月 NASA“深度撞击号”（Deep Impact）任务向坦普尔 1 号彗星（Comet 9P/Tempel 1）发送“撞击器”探测器并观测其碰撞后的结果。2014 年的一项研究表明，彗星类似于“炸冰淇淋”，即其表面由密集的晶体冰与有机化合物混合构成，而内部冰层则更冷且密度更低\(^\text{[21]}\)。

彗核表面通常干燥、多尘或岩质，表明冰层隐藏在厚度数米的表层壳体之下。彗核中含有多种有机化合物，其中可能包括甲醇、氰化氢、甲醛、乙醇、乙烷，以及可能更复杂的分子，如长链烃类与氨基酸\(^\text{[22][23]}\)。2009 年确认，NASA “星尘号”（Stardust）任务回收的彗星尘埃中发现了氨基酸甘氨酸\(^\text{[24]}\)。2011 年 8 月，一份基于 NASA 对地球陨石研究的报告提出，DNA 与 RNA 的组成单元（腺嘌呤、鸟嘌呤及相关有机分子）可能形成于小行星和彗星\(^\text{[25][26]}\)。

彗核外层表面具有极低反照率，使其成为太阳系中反射率最低的天体之一。贾托号（Giotto）探测器发现，哈雷彗星（1P/Halley）的彗核仅反射约四个百分点的入射光\(^\text{[27]}\)；而“深空一号”（Deep Space 1）发现波罗雷彗星（Comet Borrelly）的表面反射率甚至不足 3.0\%\(^\text{[27]}\)。相比之下，沥青可反射约 7\%。彗核表面的暗色物质可能由复杂的有机化合物构成。太阳加热会驱散较轻、易挥发的化合物，而留下更大的有机分子，这些物质通常极为黝黑，如焦油或原油。彗核表面的低反照率使其能够吸收驱动释气过程的热量\(^\text{[28]}\)。

已观测到的彗核半径可达 30 千米（19 英里）\(^\text{[29]}\)，但要准确测定彗核的大小并不容易\(^\text{[30]}\)。322P/SOHO 的彗核直径可能只有 100–200 米（330–660 英尺）\(^\text{[31]}\)。尽管观测仪器的灵敏度不断提高，但仍缺乏对更小彗星的探测，这使部分研究者推测直径小于 100 米（330 英尺）的彗星数量可能确实稀少\(^\text{[32]}\)。已知彗星的平均密度估计约为 $0.6 g/cm^3$（0.35 oz/cu in）\(^\text{[33]}\)。由于质量较低，彗核不会在自引力作用下成为球形，因此呈现不规则形状\(^\text{[34]}\)。
\begin{figure}[ht]
\centering
\includegraphics[width=8cm]{./figures/7e9645911c29b572.png}
\caption{彗星 81P/Wild 在光照面和阴暗面都呈现出喷流，其表面起伏鲜明，且十分干燥。} \label{fig_Comet_5}
\end{figure}
大约六个百分点的近地小行星被认为是不再发生释气作用的彗星灭绝彗核\(^\text{[35]}\)，其中包括 14827 Hypnos 与 3552 Don Quixote。

“罗塞塔号”（Rosetta）与“菲莱”（Philae）着陆器的探测结果显示，67P/丘留莫夫－格拉西缅科彗星（67P/Churyumov–Gerasimenko）的彗核不存在磁场\(^\text{[36][37]}\)，这表明磁性可能并未在早期微行星形成过程中发挥作用。此外，罗塞塔号上的 ALICE 光谱仪确定，在距离彗核约 1 千米范围内，由太阳辐射对水分子进行光电离所产生的电子，而非此前认为的来自太阳的光子，是导致彗核释放出的水与二氧化碳分子在彗发中发生分解的主要因素\(^\text{[38][39]}\)。菲莱着陆器上的仪器在彗星表面检测到至少十六种有机化合物，其中有四种（乙酰胺、丙酮、异氰酸甲酯和丙醛）是首次在彗星上被发现的\(^\text{[40][41][42]}\)。
\begin{figure}[ht]
\centering
\includegraphics[width=10cm]{./figures/8c2117a476d375ac.png}
\caption{} \label{fig_Comet_6}
\end{figure}
\subsubsection{彗发}
\begin{figure}[ht]
\centering
\includegraphics[width=8cm]{./figures/3f4fb64d5886d128.png}
\caption{Borrelly 彗星呈现喷流，但其表面不存在冰。} \label{fig_Comet_7}
\end{figure}
\begin{figure}[ht]
\centering
\includegraphics[width=8cm]{./figures/2d693476f2b0c944.png}
\caption{近日点前不久由哈勃望远镜拍摄的 ISON 彗星图像[50]} \label{fig_Comet_8}
\end{figure}
由此释放出的尘埃和气体会在彗星周围形成一个巨大而极其稀薄的大气层，称为“彗发”（coma）。太阳的辐射压与太阳风施加在彗发上的力会形成一条巨大的彗尾，且彗尾始终指向远离太阳的方向\(^\text{[51]}\)。

彗发通常由水和尘埃组成，其中水占彗星在距太阳 3 至 4 个天文单位（4.5 亿至 6 亿千米；2.8 亿至 3.7 亿英里）范围内从彗核中流出的挥发物的 90\% 之多\(^\text{[52]}\)。母体分子 $H_2O$主要通过光解作用被破坏，而光电离的贡献要小得多；与光化学过程相比，太阳风对水的破坏作用仅占很小比例\(^\text{[52]}\)。较大的尘埃颗粒会沿着彗星轨道路径残留，而较小的颗粒则被阳光压力推向远离太阳的方向进入彗尾\(^\text{[53]}\)。

虽然彗星的固体彗核通常小于 60 千米（37 英里）宽，但彗发的大小可能达到数千甚至数百万千米，有时甚至比太阳还大\(^\text{[54]}\)。例如，2007 年 10 月的一次爆发过后大约一个月，彗星 17P/Holmes 的稀薄尘埃大气层短暂地变得比太阳还大\(^\text{[55]}\)。1811 年的“大彗星”其彗发的直径大约与太阳相当\(^\text{[56]}\)。尽管彗发可以变得非常庞大，但其尺寸会在彗星经过火星轨道附近（约 1.5 AU，即 2.2 亿千米；1.4 亿英里）时缩小\(^\text{[56]}\)。在这一距离，太阳风强度足以将彗发中的气体和尘埃吹散，从而反而使彗尾变得更长\(^\text{[56]}\)。离子彗尾的长度可达到一个天文单位（1.5 亿千米）或更长\(^\text{[55]}\)。
\begin{figure}[ht]
\centering
\includegraphics[width=8cm]{./figures/f1b7d1bdbd91578d.png}
\caption{C/2006 W3（Christensen）释放碳气体（红外图像）} \label{fig_Comet_9}
\end{figure}
彗发与彗尾都由太阳照亮，并可能在彗星穿过内太阳系时变得可见；其中尘埃直接反射阳光，而气体则因电离而产生辉光\(^\text{[57]}\)。大多数彗星过于暗淡，以至于没有望远镜无法看见，但每隔数十年会有少数彗星亮到可被肉眼观测\(^\text{[58]}\)。有时彗星会经历一次巨大而突然的气体与尘埃爆发，在此期间彗发的尺寸会显著增大，持续一段时间。2007 年的霍姆斯彗星（Comet Holmes）便发生了这一现象\(^\text{[59]}\)。

1996 年，人们发现彗星能够发射 X 射线\(^\text{[60]}\)。这一发现令天文学家非常惊讶，因为 X 射线通常与极高温的天体相关。Thomas E. Cravens 在 1997 年初首次提出了解释\(^\text{[61]}\)。X 射线的产生源于彗星与太阳风之间的相互作用：当高电荷态的太阳风离子穿过彗发大气层时，它们会与彗星的原子和分子发生碰撞，并在这一过程中“夺取”一个或多个电子，这一机制称为“电荷交换”（charge exchange）。电子被太阳风离子俘获之后，会通过向基态去激发的过程发射出 X 射线和远紫外光子\(^\text{[62]}\)。
\subsubsection{弓形激波}
弓形激波是由太阳风与彗星电离层相互作用形成的，而彗星的电离层则由彗发中气体的电离所产生。随着彗星逐渐靠近太阳，释气速率不断增加，使得彗发膨胀；阳光会进一步电离彗发中的气体。当太阳风经过这一离子化的彗发区域时，弓形激波便会出现。

第一次观测是在 1980 年代与 1990 年代，当时多艘航天器先后飞掠了 21P/Giacobini–Zinner\(^\text{[63]}\)、1P/Halley\(^\text{[64]}\)与 26P/Grigg–Skjellerup\(^\text{[65]}\)彗星。结果发现，彗星的弓形激波比例如地球那样的行星弓形激波更宽且更加平缓。这些观测均发生在近日点附近，当时弓形激波已完全形成。

罗塞塔号（Rosetta）在 67P/Churyumov–Gerasimenko 彗星上观测到弓形激波的早期形成阶段，当时彗星正向太阳靠近，释气活动逐渐增强。这个尚在成长中的弓形激波被称为“婴儿弓形激波”（infant bow shock）。婴儿弓形激波是不对称的，并且相对于彗核距离而言，比完全发育的弓形激波更为宽广\(^\text{[66]}\)。
\subsubsection{彗尾}
\begin{figure}[ht]
\centering
\includegraphics[width=8cm]{./figures/ec7df3f317631c33.png}
\caption{彗星在接近太阳轨道时彗尾的典型指向方向} \label{fig_Comet_10}
\end{figure}
在外太阳系中，彗星保持冻结且处于不活跃状态，由于体积很小，从地球上极难或几乎不可能被探测到。基于哈勃空间望远镜的观测，有研究报告在柯伊伯带中统计性地探测到不活跃的彗星彗核\(^\text{[67][68]}\)，但这些探测结果曾受到质疑\(^\text{[69][70]}\)。当彗星进入内太阳系时，太阳辐射会使其中的挥发性物质升华并从彗核中逸出，同时携带尘埃。

逸出的尘埃与气体分别形成各自独立的彗尾，方向略有不同。尘埃彗尾通常沿彗星轨道附近被遗留，从而形成一条常常呈弯曲状的彗尾，称为 II 型或尘埃彗尾（dust tail）\(^\text{[57]}\)。同时，离子彗尾（ion tail），亦称 I 型彗尾，由气体组成，始终直接指向远离太阳的方向，因为气体比尘埃更强烈地受到太阳风作用，并沿磁力线运动，而非遵循轨道轨迹\(^\text{[71]}\)。在某些情况下——例如当地球经过彗星的轨道平面时——可以看到一条反尾（antitail），其方向与离子彗尾和尘埃彗尾相反\(^\text{[72]}\)。
\begin{figure}[ht]
\centering
\includegraphics[width=8cm]{./figures/1dacb212df9b211e.png}
\caption{彗星示意图：展示由太阳风形成的尘埃轨迹、尘埃彗尾与离子气体彗尾} \label{fig_Comet_11}
\end{figure}
对反尾（antitail）的观测对太阳风的发现起到了重要作用\(^\text{[73]}\)。离子彗尾是由于彗发中粒子被太阳紫外辐射电离而形成的。一旦粒子被电离，它们便获得净正电荷，从而在彗星周围产生一个“感应磁层”（induced magnetosphere）。彗星及其感应磁场对向外流动的太阳风粒子构成阻碍。由于彗星与太阳风之间的相对轨道速度为超音速，彗星在太阳风流向的上游会形成一个弓形激波。在这一弓形激波中，彗星离子（称为“拾取离子” pick-up ions）大量聚集，并使太阳磁场被等离子体“负载”，使磁力线在彗星周围“披覆”（drape），从而形成离子彗尾\(^\text{[74]}\)。

若离子彗尾的负载程度足够大，磁力线会被压缩到某一距离处发生磁重联（magnetic reconnection）。这会导致一次“彗尾断裂事件”（tail disconnection event）\(^\text{[74]}\)。这一现象已多次被观测到，其中一个著名事件发生在 2007 年 4 月 20 日：当恩克彗星（Encke’s Comet）穿过一次日冕物质抛射（CME）时，其离子彗尾完全被切断。该事件由 STEREO 空间探测器记录\(^\text{[75]}\)。

2013 年，欧洲航天局（ESA）的科学家报告称，金星的电离层会向外流出，其行为与类似条件下彗星的离子彗尾相似\(^\text{[76][77]}\)。
\subsubsection{喷发}
\begin{figure}[ht]
\centering
\includegraphics[width=8cm]{./figures/e97b94991a13b0fc.png}
\caption{103P/Hartley 的气体与冰雪喷流} \label{fig_Comet_12}
\end{figure}
不均匀的加热会使新生成的气体从彗核表面的薄弱区域喷出，类似于间歇泉\(^\text{[78]}\)。这些气体和尘埃的喷流能够使彗核旋转，甚至将其分裂\(^\text{[78]}\)。2010 年的研究显示，干冰（冻结的二氧化碳）的升华能够为从彗核喷出的物质喷流提供动力\(^\text{[79]}\)。对 Hartley 2 彗星的红外成像显示了此类喷流从其表面喷出，并携带尘埃颗粒进入彗发\(^\text{[80]}\)。
\subsection{轨道特征}
大多数彗星是太阳系内的小天体，其轨道为细长的椭圆轨道，轨道的一部分将它们带到靠近太阳的区域，而其余时间则位于太阳系的遥远外侧\(^\text{[81]}\)。彗星通常根据其轨道周期长短进行分类：周期越长，其轨道椭圆的离心率越大。
\subsubsection{短周期}
周期彗星或短周期彗星通常被定义为那些轨道周期小于 200 年的彗星\(^\text{[82]}\)。它们通常在黄道面内、并以与行星相同的方向运行\(^\text{[83]}\)。其轨道在远日点时往往延伸至外行星（木星及更远）的区域；例如，哈雷彗星的远日点略微超出海王星轨道。远日点接近某一主要行星轨道的彗星被称为该行星的“族”（family）\(^\text{[84]}\)。一般认为，这类彗星族是由行星将原本属于长周期彗星的天体捕获至较短轨道而形成的\(^\text{[85]}\)。

在短轨道周期的极端情况下，恩克彗星（Encke's Comet）的轨道甚至未到达木星轨道，被称为恩克型彗星（Encke-type comet）。轨道周期小于 20 年、且轨道倾角对黄道面不超过 30 度的短周期彗星被称为传统木星族彗星（Jupiter-family comets, JFCs）\(^\text{[86][87]}\)。类似哈雷彗星、轨道周期介于 20 至 200 年之间、且轨道倾角范围从 0 度一直到超过 90 度的，则称为哈雷型彗星（Halley-type comets, HTCs）\(^\text{[88][89]}\)。截至 2025 年 7 月，已知恩克型彗星共有 74 颗（其中 6 颗被归类为近地天体 NEOs），HTCs 有 109 颗（其中 36 颗为 NEOs），而 JFCs 有 815 颗（其中 153 颗为 NEOs）\(^\text{[90]}\)。

近期发现的主带彗星（main-belt comets）构成一个独立的类别，其轨道更为接近圆形，并位于小行星带内部运行\(^\text{[91][92]}\)。

由于其椭圆轨道经常使它们接近巨行星，彗星会受到进一步的引力扰动\(^\text{[93]}\)。短周期彗星的远日点往往与某一巨行星的半长轴相一致，其中木星族彗星（JFCs）是最大的一类\(^\text{[87]}\)。来自奥尔特云的彗星在经历与巨行星的近距离遭遇后，其轨道常会受到强烈的引力影响。木星因其质量超过其他所有行星的总和，是引力扰动最显著的来源。这些扰动能够将长周期彗星偏转至较短周期的轨道上\(^\text{[94][95]}\)。

基于轨道特征，短周期彗星通常被认为起源于半人马小行星（centaurs）以及柯伊伯带／散射盘（Kuiper belt/scattered disc）\(^\text{[96]}\) —— 一个位于海王星外侧区域的盘状小天体群；而长周期彗星的来源则被认为是更遥远的球形奥尔特云（Oort cloud），其名称来自假设其存在的荷兰天文学家 Jan Hendrik Oort\(^\text{[97]}\)。在这些遥远区域中，巨量彗星状天体被认为以近圆轨道绕太阳运行。有时外行星（对柯伊伯带天体而言）或邻近恒星（对奥尔特云天体而言）的引力会将其中某个天体抛入一条椭圆轨道，使其向太阳方向运动，从而形成可见彗星。与周期彗星的回归不同，这类新彗星的出现是不可预测的，因为它们的轨道此前从未被建立\(^\text{[98]}\)。当地被抛入太阳的轨道并不断向太阳落入时，其质量会持续被剥离，而这种剥离对彗星的寿命影响巨大：剥离越严重，寿命越短，反之亦然\(^\text{[99]}\)。
\subsubsection{长周期}
\begin{figure}[ht]
\centering
\includegraphics[width=8cm]{./figures/aff8f46e648a6626.png}
\caption{科霍特克彗星（红色）与地球（蓝色）的轨道示意图，展示其轨道的高离心率以及彗星在接近太阳时的高速运动。} \label{fig_Comet_13}
\end{figure}
长周期彗星具有高度偏心的轨道，其周期范围从 200 年延伸至数千年、甚至数百万年\(^\text{[100]}\)。在近日点附近出现偏心率大于 1 的情况并不必然意味着彗星将离开太阳系\(^\text{[101]}\)。例如，麦克诺特彗星（Comet McNaught）在 2007 年 1 月经过近日点时，其日心瞬时轨道偏心率为 1.000019，但其真实轨道仍被太阳束缚，轨道周期约为 92,600 年，因为偏心率在其远离太阳后会下降至 1 以下。要得到长周期彗星的未来轨道，必须在其离开行星区域后计算其瞬时轨道，并且应相对于太阳系质心来计算。根据定义，长周期彗星保持被太阳引力束缚；那些因与主要行星近距离掠过而被抛出太阳系的彗星，不再被认为具有“周期”。长周期彗星的远日点远远超出外行星，其轨道平面也不必与黄道面接近。诸如 C/1999 F1 与 C/2017 T2（PANSTARRS）等长周期彗星，其远日点距离可以接近 70,000 AU（0.34 pc；1.1 光年），估计轨道周期约为 600 万年。

单次显现或非周期彗星在性质上与长周期彗星类似，因为它们在内太阳系近日点附近具有近似抛物线或略呈双曲线的轨迹\(^\text{[100]}\)。然而，巨行星的引力扰动会改变它们的轨道。单次显现彗星的瞬时轨道为双曲线或抛物线，使其在仅经过一次太阳后便永久离开太阳系\(^\text{[102]}\)。太阳的希尔球（Hill sphere）具有约 230,000 AU（1.1 pc；3.6 光年）的不稳定最大边界\(^\text{[103]}\)。只有数百颗彗星在近日点附近被观测到具有双曲线轨道（e > 1）\(^\text{[104]}\)；基于日心的、未受扰动的双体拟合轨道表明，它们可能会逃离太阳系。

截至 2025 年，已有三颗偏心率显著大于 1 的天体被发现：1I/ʻOumuamua、2I/Borisov 与 3I/ATLAS，它们均被认为起源于太阳系之外。ʻOumuamua 的偏心率约为 1.2，在其于 2017 年 10 月穿过内太阳系期间并未显示出彗星活动的光学迹象，但其轨迹的变化——暗示存在释气作用——表明其很可能是一颗彗星\(^\text{[105]}\)。相比之下，2I/Borisov 的偏心率约为 3.36，已被观测到具有典型的彗发结构，被认为是首次被探测到的星际彗星\(^\text{[106][107]}\)。3I/ATLAS 的偏心率约为 6.1，同样具有彗发，说明它也是一颗彗星。C/1980 E1 在 1982 年近日点之前的轨道周期约为 710 万年，但其在 1980 年与木星的一次遭遇使其加速，赋予它目前已知太阳系内彗星中最大的偏心率（1.057），且观测弧度相对可靠\(^\text{[108]}\)。预计不会再返回内太阳系的彗星包括 C/1980 E1、C/2000 U5、C/2001 Q4（NEAT）、C/2009 R1、C/1956 R1 以及 C/2007 F1（LONEOS）。

部分权威资料将“周期彗星”（periodic comet）用于指代所有轨道周期性的彗星（即所有短周期与长周期彗星）\(^\text{[109]}\)，而另一些资料则仅将其限定为短周期彗星\(^\text{[100]}\)。类似地，尽管“非周期彗星”（non-periodic comet）从字面上等同于“单次显现彗星”，但在另一种用法中，它被用于指代所有非上述狭义“周期彗星”的彗星（即轨道周期超过 200 年的所有彗星）。

早期观测确实发现了一些真正的双曲线（即非周期）轨迹，但数量不超过木星引力扰动能够解释的范围。来自星际空间的彗星其速度量级与太阳附近恒星的相对速度相似（每秒数十公里）。当这类天体进入太阳系时，它们具有正的轨道比能，从而具有正的无穷远速度（$v_{\infty}$），并呈现明显的双曲线轨迹。粗略估计，在木星轨道范围内，每世纪可能出现约四颗星际双曲线彗星，误差可达一个至两个数量级\(^\text{[110]}\)。
\begin{figure}[ht]
\centering
\includegraphics[width=14.25cm]{./figures/b8f95819756446e8.png}
\caption{} \label{fig_Comet_14}
\end{figure}
\subsubsection{奥尔特云与希尔斯云}
\begin{figure}[ht]
\centering
\includegraphics[width=14.25cm]{./figures/8a76749915cbfb85.png}
\caption{被认为包围着太阳系的奥尔特云示意图，并展示了柯伊伯带和小行星带以作比较。} \label{fig_Comet_15}
\end{figure}
奥尔特云被认为占据着一片广阔的空间，其内边界大约从距离太阳 2,000 至 5,000 AU（0.03–0.08 光年）\(^\text{[112]}\)开始，一直到约 50,000 AU（0.79 光年）\(^\text{[88]}\)。这一云团包围着从太阳系中部——太阳本身——一直到柯伊伯带外缘的所有天体。奥尔特云由适合形成天体的可行物质组成。太阳系行星之所以存在，是因为这些由太阳引力聚集并凝结的微行星（形成行星时遗留下来的固体碎块）。这些被困微行星形成的偏心性正是奥尔特云存在的原因\(^\text{[113]}\)。一些估计将其外缘置于 100,000 至 200,000 AU（1.58–3.16 光年）之间\(^\text{[112]}\)。该区域可进一步分为距离太阳 20,000–50,000 AU（0.32–0.79 光年）的球状外奥尔特云，以及距离太阳 2,000–20,000 AU（0.03–0.32 光年）的甜甜圈状内云，即希尔斯云（Hills cloud）\(^\text{[114]}\)。外奥尔特云与太阳的引力束缚很弱，是落入海王星轨道以内的长周期彗星（以及可能的哈雷型彗星）的来源\(^\text{[88]}\)。内奥尔特云亦被称为希尔斯云，以 1981 年提出其存在的 Jack G. Hills 命名\(^\text{[115]}\)。模型预测内云中包含的彗核数量应为外云的数十倍至数百倍\(^\text{[115][116][117]}\)；其被视为新彗星的潜在来源，为随着时间流逝而逐渐变得稀薄的外奥尔特云提供补给。希尔斯云解释了奥尔特云在数十亿年后仍然存在的原因\(^\text{[118]}\)。
\subsubsection{系外彗星}
太阳系之外的系外彗星已被探测到，并可能在银河系中普遍存在\(^\text{[119]}\)。首个被探测到的系外彗星系统是 1987 年发现的织女增二（Beta Pictoris）系统，它是一颗极年轻的 A 型主序星\(^\text{[120][121]}\)。截至 2013 年，共识别出 11 个此类系外彗星系统，其探测方式依赖于彗星在靠近其母恒星时释放的大量气体所造成的吸收光谱特征\(^\text{[119][120]}\)。开普勒空间望远镜在十年间负责搜索太阳系外的行星与其他天体。首例凌日系外彗星于 2018 年 2 月被一个由专业天文学家与公民科学家组成的小组发现，他们在开普勒的光变曲线中识别出了这些信号\(^\text{[122][123]}\)。在开普勒望远镜于 2018 年 10 月退役后，其任务由 TESS 望远镜接替。自 TESS 发射以来，天文学家利用其光变曲线在织女增二周围观测到彗星的凌日现象\(^\text{[124][125]}\)。

随着 TESS 的投入使用，天文学家能够更好地利用光谱方法区分系外彗星。新行星通常通过白光光变法探测，当行星遮挡母恒星时，会在图表读数中出现对称的下降。然而，对这些光变曲线的进一步分析显示，这些不对称下降的模式实际上是由彗尾或数百条彗尾引起的\(^\text{[126]}\)。
\subsection{彗星的影响}
\begin{figure}[ht]
\centering
\includegraphics[width=8cm]{./figures/b8829ab21e3c6012.png}
\caption{英仙座流星示意图} \label{fig_Comet_16}
\end{figure}
\subsubsection{彗星与流星雨的关联}
当彗星在接近太阳的过程中受热时，其冰质组分的释气作用会释放出固体碎屑，这些碎屑太大，辐射压与太阳风无法将其完全吹散\(^\text{[127]}\)。如果地球的轨道恰好穿过这条主要由细小岩质颗粒组成的碎屑轨迹，那么当地球经过其中时就可能出现流星雨。碎屑轨迹越密集，产生的流星雨越短暂但越强烈；而较稀疏的碎屑轨迹则会形成时间较长但强度较低的流星雨。通常而言，碎屑轨迹的密度与母彗星释放这些物质的时间有直接关系\(^\text{[128][129]}\)。例如，英仙座流星雨每年 8 月 9 日至 13 日出现，此时地球通过斯威夫特－塔特尔彗星（Comet Swift–Tuttle）的轨道。哈雷彗星则是 10 月出现的猎户座流星雨的来源\(^\text{[130][131]}\)。
\subsubsection{彗星与生命的影响}
在早期地球形成阶段，许多彗星和小行星与地球发生过碰撞。许多科学家认为，大约 40 亿年前彗星轰击年轻地球，为地球海洋带来了如今巨量的水，或者至少带来了其中相当大的一部分。然而，也有研究对这一观点提出质疑\(^\text{[132]}\)。在彗星中探测到的有机分子（包括多环芳烃）\(^\text{[21]}\)的显著含量，引发了关于彗星或陨石可能将生命前体——甚至生命本身——带到地球的推测\(^\text{[133]}\)。2013 年有研究提出，像彗星这样岩石表面与冰表面发生碰撞的事件具有通过冲击合成生成构成蛋白质的氨基酸的潜力\(^\text{[134]}\)。彗星进入大气层的高速运动，以及初次撞击产生的巨大能量，使得小分子能够凝结并形成更大的大分子，这些大分子为生命的出现奠定了基础\(^\text{[135]}\)。2015 年，科学家在 67P 彗星的释气物质中发现了大量的分子氧，这暗示该分子可能比原先认为的更为常见，从而使其作为生命迹象的意义降低\(^\text{[136]}\)。

人们推测，在长时间尺度上，彗星撞击可能向月球输送了大量水，其中一部分可能以月球冰的形式保留下来\(^\text{[137]}\)。彗星和流星体撞击也被认为是玻陨石（tektites）和澳洲玻陨石（australites）存在的原因\(^\text{[138]}\)。
\subsubsection{对彗星的恐惧}
在 1200 至 1650 年间，欧洲人对彗星的恐惧达到高峰，将其视为上帝的行为或灾难即将来临的征兆\(^\text{[139]}\)。例如，在 1618 年“大彗星”出现后的第二年，Gotthard Arthusius 发表了一份小册子，声称该彗星预示着世界末日临近\(^\text{[140]}\)。他列出了十页与彗星相关的灾难，包括“地震、洪水、河道改变、冰雹、炎热干燥的天气、歉收、瘟疫、战争、背叛与物价飞涨”等\(^\text{[139]}\)。

到 1700 年，大多数学者已经认识到这些灾难无论是否出现彗星都会发生。然而，威廉·惠斯顿（William Whiston）在 1711 年利用爱德蒙·哈雷记录的彗星观测数据写道，1680 年的大彗星具有 574 年的周期性，并通过向地球倾泻洪水的方式导致了《创世纪》中的全球性大洪水。他的观点再次激发了人们对彗星的恐惧，这一次不是将其视为灾难的征兆，而是视为地球的直接威胁\(^\text{[139]}\)。1910 年光谱分析在哈雷彗星彗尾中发现了有毒气体氰（cyanogen）\(^\text{[141]}\)，导致公众恐慌性购买防毒面具以及各种伪科学的“抗彗星药片”和“抗彗星雨伞”\(^\text{[142]}\)。
\subsection{彗星的命运}
\subsubsection{离开（被抛出）太阳系}
如果一颗彗星的速度足够快，它可能会离开太阳系。这样的彗星遵循双曲线的开放轨道，因此被称为双曲线彗星。已知的太阳系内双曲线彗星只有在与太阳系中其他天体（如木星）相互作用后才会被抛出\(^\text{[143]}\)。一个例子是 C/1980 E1 彗星，它在 1980 年与木星的一次近距离掠过后，其原本约 710 万年的绕日轨道被改变为双曲线轨道\(^\text{[144]}\)。而诸如 1I/ʻOumuamua、2I/Borisov 和 3I/ATLAS 等星际彗星从未绕太阳运行，因此不需要第三天体的相互作用即可离开太阳系。
\subsubsection{灭绝}
木星族彗星（JFCs）与长周期彗星表现出非常不同的亮度衰减规律。木星族彗星的活跃寿命约为 10,000 年，或大约 1,000 个轨道周期，而长周期彗星的衰减速度要快得多。只有 10\% 的长周期彗星能在多于 50 次的小近日点掠过中存活，而只有 1\% 能存活超过 2,000 次掠过\(^\text{[35]}\)。最终，彗核中的大部分挥发性物质都会蒸发，彗星会变成一个小型、黑暗、惰性的岩石或碎块，其外观可能类似小行星\(^\text{[145]}\)。一些椭圆轨道的小行星如今被识别为灭绝彗星\(^\text{[146][147][148][149]}\)。大约六个百分点的近地小行星被认为是灭绝的彗星彗核\(^\text{[35]}\)。
\subsubsection{破裂与碰撞}
某些彗星的彗核可能十分脆弱，这一结论得到了彗星分裂现象的支持\(^\text{[150]}\)。一次重大彗星破裂事件是 1993 年发现的舒梅克－列维 9 号彗星（Comet Shoemaker–Levy 9）。其在 1992 年 7 月与木星的一次近距离遭遇中被撕裂成多个碎片，并在 1994 年 7 月的六天时间里，这些碎片陆续坠入木星大气层——这是天文学家首次观测到太阳系内两个天体之间的碰撞\(^\text{[151][152]}\)。其他发生分裂的彗星包括 1846 年的 3D/Biela 以及 1995 至 2006 年间的 73P/Schwassmann–Wachmann\(^\text{[153]}\)。希腊历史学家埃福鲁斯（Ephorus）报告称，早在公元前 372–373 年的冬季就曾有彗星分裂\(^\text{[154]}\)。人们认为彗星可能因热应力、内部气体压力或撞击而发生分裂\(^\text{[155]}\)。

42P/Neujmin 和 53P/Van Biesbroeck 彗星似乎是同一母体彗星的碎片。数值积分显示，这两颗彗星在 1850 年 1 月曾非常接近木星，且在 1850 年之前，它们的轨道几乎完全相同\(^\text{[156]}\)。另一个由分裂事件产生的彗星族是 Liller 彗星家族，包括 C/1988 A1（Liller）、C/1996 Q1（Tabur）、C/2015 F3（SWAN）、C/2019 Y1（ATLAS）以及 C/2023 V5（Leonard）\(^\text{[157][158]}\)。

一些彗星在近日点通过期间发生了破裂，包括著名的大彗星 West 和 Ikeya–Seki。比埃拉彗星（Biela's Comet）是一个典型例子：它在 1846 年穿过近日点时分裂为两个部分。这两个部分在 1852 年仍被分开观测到，但此后再也未被看到。相反，在 1872 年和 1885 年（彗星本应出现的时段）出现了壮观的流星雨。每年 11 月出现的小型流星雨——仙女座流星雨（Andromedids）——便是在地球穿过比埃拉彗星轨道时产生的\(^\text{[159]}\)。

有些彗星则以更加戏剧性的方式终结——要么坠入太阳\(^\text{[160]}\)，要么与行星或其他天体发生碰撞。在早期太阳系中，彗星与行星或卫星的碰撞十分常见：例如月球上的许多陨石坑可能是彗星造成的。最近一次彗星与行星的碰撞发生在 1994 年 7 月，当时舒梅克－列维 9 号彗星分裂成碎片并与木星相撞\(^\text{[161]}\)。
\begin{figure}[ht]
\centering
\includegraphics[width=14.25cm]{./figures/fad311c089faff94.png}
\caption{} \label{fig_Comet_17}
\end{figure}
\subsection{命名法}
\begin{figure}[ht]
\centering
\includegraphics[width=8cm]{./figures/fcac94c87ed8f3ed.png}
\caption{1910 年的哈雷彗星} \label{fig_Comet_18}
\end{figure}
过去两个世纪中，彗星的命名遵循过几种不同的惯例。在 20 世纪初之前，大多数彗星以其出现的年份命名，有时也会为特别明亮的彗星加上修饰词；因此出现了诸如“1680 年大彗星”、“1882 年大彗星”以及“1910 年一月大彗星”等名称。

在爱德蒙·哈雷（Edmond Halley）证明 1531 年、1607 年和 1682 年出现的彗星其实是同一天体，并通过计算其轨道成功预测其将在 1759 年回归之后，该彗星便被称为哈雷彗星（Halley's Comet）\(^\text{[163]}\)。类似地，第二与第三颗已知的周期彗星——恩克彗星（Encke's Comet）\(^\text{[164]}\)和比埃拉彗星（Biela's Comet）\(^\text{[165]}\)——以计算其轨道的天文学家而非最初的发现者命名。后来，周期彗星通常以其发现者命名，但仅出现一次的彗星仍通常以其出现的年份命名\(^\text{[166]}\)。

20 世纪初，人们开始普遍采用以发现者命名彗星的惯例，直到今日仍是如此。一颗彗星可以以其发现者命名，也可以以发现它的仪器或科学计划命名\(^\text{[166]}\)。例如，2019 年，天文学家 Gennadiy Borisov 观测到一颗似乎起源于太阳系之外的彗星；该彗星被命名为 2I/Borisov，以纪念他的发现\(^\text{[167]}\)。
\subsection{研究历史}
\subsubsection{早期观测与思想}
\begin{figure}[ht]
\centering
\includegraphics[width=8cm]{./figures/cfd768f94e0b971f.png}
\caption{占星手稿细节，墨绘于丝帛，西汉公元前 2 世纪，出土自马王堆墓。此页描述并绘制了七颗彗星，在整部文献中共记录了 29 颗（参见：中国历史彗星观测）[171]。} \label{fig_Comet_19}
\end{figure}
从古代文献（如中国甲骨文）可知，人类在数千年前就已经注意到彗星的存在\(^\text{[168]}\)。直到十六世纪之前，彗星通常被视为不祥之兆，被认为预示国王或贵族之死、灾难将临，甚至被解读为天界生物对地上居民的攻击\(^\text{[169][170]}\)。

亚里士多德（Aristotle，公元前 384–322 年）是已知第一位运用多种理论与观测事实，建立连贯且结构化的彗星宇宙学理论的科学家。他认为彗星是大气层中的现象，因为它们可以在黄道带之外出现，并且在数日内亮度会发生变化。亚里士多德的彗星理论源于他的观测以及其宇宙论：宇宙中的一切都按照特定的结构方式排列\(^\text{[172]}\)。在这种结构中，天界与地界被明确区分，而彗星被严格视为属于地界的现象。根据亚里士多德的观点，彗星必定位于月球轨道之内，并与天界存在明显分隔。同样在公元前 4 世纪，米恩多斯的阿波罗尼乌斯（Apollonius of Myndus）支持彗星像行星一样运动的观点\(^\text{[173]}\)。尽管中世纪期间已有多项发现质疑其理论的部分内容，但亚里士多德的彗星学说仍长期被广泛接受\(^\text{[174]}\)。

在公元 1 世纪，小塞涅卡（Seneca the Younger）开始质疑亚里士多德关于彗星的逻辑。他指出，由于彗星运动规律且不受风的影响，它们不可能是大气现象\(^\text{[175]}\)，并且它们比在天空中一闪而过的短暂表象更加持久。\(^\text{[a]}\)他指出只有彗尾才是透明的、云状的，并论证没有理由将彗星的轨道限制在黄道带之内\(^\text{[175]}\)。在批评米恩多斯的阿波罗尼乌斯时，塞涅卡写道：“彗星穿过宇宙的高层区域，最终在其轨道的最低点才变得可见。”\(^\text{[176]}\)虽然塞涅卡本人并未建立一个完整的彗星理论\(^\text{[177]}\)，但他的论证在 16、17 世纪引发了对亚里士多德观点的大量批评与讨论\(^\text{[174][b]}\)。

在公元 1 世纪，老普林尼（Pliny the Elder）认为彗星与政治动荡和死亡有关\(^\text{[179]}\)。普林尼将彗星视为“类人”的天象，常以“长发”或“长须”等词形容其彗尾\(^\text{[180]}\)。他根据彗星颜色与形状的分类体系被沿用数个世纪\(^\text{[181]}\)。

在印度，至公元 6 世纪时，天文学家相信彗星是周期性再现的天象。6 世纪的天文学家 Varāhamihira 与 Bhadrabahu 持这种观点，而 10 世纪的天文学家 Bhaṭṭotpala 列出了若干彗星的名称及其推测周期，但这些数据的计算方式与准确性均不得而知\(^\text{[182][183]}\)。
\begin{figure}[ht]
\centering
\includegraphics[width=8cm]{./figures/60c06edf50f8a5c8.png}
\caption{哈雷彗星于 1066 年出现在黑斯廷斯战役之前，并被描绘在《拜厄克斯挂毯》中。} \label{fig_Comet_20}
\end{figure}
有一种说法认为，一位阿拉伯学者在 1258 年记录了数次彗星（或某类彗星）的重复出现，尽管尚不清楚他是否将其视为同一颗周期彗星，但该彗星的周期可能约为 63 年\(^\text{[184]}\)。

1301 年，意大利画家乔托（Giotto）成为第一位以准确且具解剖学特征的方式描绘彗星的人。在他的作品《三王来朝》（Adoration of the Magi）中，乔托以哈雷彗星来替代伯利恒之星的描绘，其准确程度直到 19 世纪才被超越，而真正被超越则是在摄影术发明之后\(^\text{[185]}\)。

尽管现代科学天文学已开始扎根，但在 15 世纪之前，占星式的彗星解读仍占主导地位。彗星依旧被视为灾难的预兆，如《卢塞恩编年史》（Luzerner Schilling）中的描述以及教皇卡利克斯图斯三世（Pope Callixtus III）的警告所示\(^\text{[185]}\)。1578 年，德国路德宗主教 Andreas Celichius 将彗星定义为“人类罪恶的浓烟……被至高天庭审判者炙热而燃烧的怒火所点燃”。翌年，Andreas Dudith 则指出：“如果彗星是由凡人的罪孽引起，那么它们就永远不会从天空中消失。”\(^\text{[186]}\)
\subsubsection{现代天文学}
1456 年，人们曾对哈雷彗星进行粗略的视差测量尝试，但结果是错误的\(^\text{[187]}\)。雷吉蒙塔努斯（Regiomontanus）是第一位试图通过观测 1472 年大彗星来计算日平行（diurnal parallax）的人。他的预测并不十分准确，但其目的在于尝试估算彗星与地球之间的距离\(^\text{[181]}\)。
\begin{figure}[ht]
\centering
\includegraphics[width=8cm]{./figures/dba9ef1f7901154b.png}
\caption{第谷·布拉赫在其笔记本中绘制的 1577 年大彗星观测草图} \label{fig_Comet_21}
\end{figure}
在 16 世纪，第谷·布拉赫（Tycho Brahe）和迈克尔·梅斯特林（Michael Maestlin）通过测量 1577 年大彗星的视差，证明彗星必定存在于地球大气层之外\(^\text{[188]}\)。在当时测量精度的条件下，这意味着彗星的距离至少是地月距离的四倍\(^\text{[189][190]}\)。基于 1664 年的观测，乔万尼·博雷利（Giovanni Borelli）记录了他所观测彗星的经度与纬度，并提出彗星轨道可能呈抛物线形状\(^\text{[191]}\)。尽管伽利略（Galileo Galilei）是一位技艺精湛的天文学家，但他在 1623 年的著作《试金石》（The Assayer）中否定了第谷的彗星视差理论，声称彗星可能只是视觉幻象，而他本人几乎没有进行过相关观测\(^\text{[181]}\)。1625 年，梅斯特林的学生约翰内斯·开普勒（Johannes Kepler）确认第谷关于彗星视差的观点是正确的\(^\text{[181]}\)。此外，数学家雅各布·伯努利（Jacob Bernoulli）于 1682 年发表了一部关于彗星的论文。

在近代早期，彗星也因其在医学中的占星意义而被研究。当时许多治疗者认为医学与天文学密不可分，他们利用自己对彗星及其他占星征兆的理解来诊断与治疗病人\(^\text{[192]}\)。

艾萨克·牛顿在 1687 年的《自然哲学的数学原理》（Principia Mathematica）中证明：若一个物体在平方反比引力作用下运动，其轨道必定为圆锥曲线之一。他还展示了如何将一颗彗星在天空中的路径拟合为抛物线轨道，以 1680 年的大彗星为例\(^\text{[193]}\)。他将彗星描述为在斜轨道上运动的致密而耐久的固体天体，并认为彗尾是由彗核释放出的稀薄蒸气，在阳光下被点燃或加热。他怀疑彗星是空气中维持生命成分的来源\(^\text{[194]}\)。他指出彗星通常出现在太阳附近，因此极有可能绕太阳运行\(^\text{[175]}\)。关于彗星的亮度，他写道：“彗星因反射太阳光而发光”，其彗尾则由“阳光照亮由（彗发）升起的烟雾”构成\(^\text{[175]}\)。
\begin{figure}[ht]
\centering
\includegraphics[width=8cm]{./figures/2687c643396781aa.png}
\caption{} \label{fig_Comet_22}
\end{figure}
1705 年，爱德蒙·哈雷（Edmond Halley，1656–1742）将牛顿的方法应用于 1337 年至 1698 年间出现的 23 次彗星现象。他注意到其中的三次——1531 年、1607 年和 1682 年的彗星——具有非常相似的轨道要素，并进一步能够将它们轨道间的轻微差异解释为由木星与土星引力扰动所造成的结果。确信这三次出现均为同一颗彗星的三次回归后，他预测该彗星将在 1758–59 年再次出现\(^\text{[195]}\)。

哈雷所预测的回归日期随后被三位法国数学家——亚历克西·克莱罗（Alexis Clairaut）、约瑟夫·拉朗德（Joseph Lalande）和妮可-蕾娜·勒波特（Nicole-Reine Lepaute）——进一步修正，他们将彗星 1759 年近日点的日期预测到了一个月以内的精度\(^\text{[196][197]}\)。当彗星如期回归后，它便被命名为哈雷彗星（Halley's Comet）\(^\text{[198]}\)。

早在 18 世纪，就已有科学家对彗星的物理组成提出了正确的假设。1755 年，伊曼努尔·康德（Immanuel Kant）在其《通用自然史》中假设彗星是由“原始物质”在已知行星之外凝聚而成，这些物质受引力影响“微弱移动”，从而以任意倾角运行，并在靠近日心点时部分被太阳的热量汽化\(^\text{[200]}\)。1836 年，德国数学家弗里德里希·威廉·贝塞尔（Friedrich Wilhelm Bessel）在观测 1835 年哈雷彗星出现时其彗发中的气体喷流后，提出蒸发物质的喷射力足以显著改变彗星轨道，并认为恩克彗星的非引力运动正是由这一现象造成的\(^\text{[201]}\)。

19 世纪，帕多瓦天文台（Astronomical Observatory of Padova）成为彗星观测研究的中心。在乔瓦尼·桑蒂尼（Giovanni Santini，1787–1877）以及后来的朱塞佩·洛伦佐尼（Giuseppe Lorenzoni，1843–1914）的领导下，该天文台致力于经典天文学，特别是新彗星与行星轨道的计算，目标是编制包含近一万颗恒星的星表。位于意大利北部的地理位置使该台的观测在建立若干重要的大地测量、地理与天文参数方面发挥了关键作用，例如米兰与帕多瓦之间的经度差，以及帕多瓦至菲乌梅的经度差\(^\text{[202]}\)。天文台内部的通信，尤其是桑蒂尼与另一位天文学家朱塞佩·托阿尔多（Giuseppe Toaldo）之间的往来信件，多次提及彗星与行星轨道观测的重要性\(^\text{[203]}\)。

1950 年，弗雷德·劳伦斯·惠普尔（Fred Lawrence Whipple）提出彗星并非含少量冰的岩质天体，而是含少量岩石与尘埃的冰质天体\(^\text{[204]}\)。这一“肮脏雪球”（dirty snowball）模型很快被接受，并在 1986 年一系列飞越哈雷彗星彗发的航天器（包括欧洲航天局的 Giotto 探测器及苏联的 Vega 1 与 Vega 2）所拍摄的彗核图像与喷流观测中得到了支持\(^\text{[205]}\)。

2014 年 1 月 22 日，欧洲航天局（ESA）的科学家首次明确探测到小行星带中最大的天体——矮行星谷神星（Ceres）上的水汽\(^\text{[206]}\)。这一探测由赫歇尔空间天文台（Herschel Space Observatory）的远红外能力完成\(^\text{[207]}\)。这一发现出乎意料，因为通常认为会产生喷流与羽状物的是彗星，而非小行星。一位科学家表示：“彗星与小行星之间的界线正在变得越来越模糊。”\(^\text{[207]}\)2014 年 8 月 11 日，天文学家首次使用阿塔卡马大型毫米/亚毫米阵列（ALMA）发布研究，展示了 C/2012 F6（Lemmon）和 C/2012 S1（ISON）彗星彗发中 HCN、HNC、H₂CO 与尘埃的分布情况\(^\text{[208][209]}\)。
\subsubsection{航天器任务}  
\begin{itemize}
\item “哈雷舰队”（Halley Armada）指 1980 年代在哈雷彗星近日点期间，对其进行访问和／或观测的一系列航天器任务。原计划在 1986 年由航天飞机“挑战者号”（Space Shuttle Challenger）执行对哈雷彗星的研究任务，但该航天飞机在发射后不久爆炸。
\item \textbf{Deep Impact.} 关于彗星内部冰的含量，学界仍存在争论。2001 年，“深空一号”（Deep Space 1）探测器获得了博雷利彗星（Comet Borrelly）表面的高分辨率图像。结果发现博雷利彗星的表面干燥且温度较高，介于 26 至 71 °C（79–160 °F），并且极其黑暗，暗示其冰层已被太阳加热与老化过程移除，或被覆盖其表面的类烟灰物质所掩藏\(^\text{[210]}\)。2005 年 7 月，“深度撞击”（Deep Impact）探测器在坦普尔 1 号彗星（Comet Tempel 1）上制造了一个撞击坑以研究其内部结构。结果表明，彗星的大部分水冰位于表面之下，这些储库为形成坦普尔 1 彗发的水汽喷流提供来源\(^\text{[211]}\)。该任务后来更名为 EPOXI，并于 2010 年 11 月 4 日飞掠哈特利 2 号彗星（Comet Hartley 2）。
\item \textbf{Ulysses.} 2007 年，“尤利西斯号”（Ulysses）探测器意外穿过 2006 年发现的 C/2006 P1（McNaught）彗尾。尤利西斯号于 1990 年发射，其原始使命是绕太阳运行并在各个纬度进行研究
\item \textbf{Stardust.} “星尘号”（Stardust）任务的数据表明，从维尔德 2 号彗星（Wild 2）彗尾回收的物质是晶质的，并且只能在超过 1,000 °C（1,830 °F）的“火中诞生”\(^\text{[212][213]}\)。虽然彗星形成于外太阳系，但太阳系早期的径向物质混合过程被认为将物质重新分布到了整个原行星盘中\(^\text{[214]}\)。因此，彗星中包含形成于早期太阳系高温内区的晶质颗粒。这一点既体现在彗星光谱中，也体现在样本返回任务中。更近期的分析显示，回收的物质表明“彗尘与小行星物质相似”\(^\text{[215]}\)。这些新结果迫使科学家重新思考彗星的本质及其与小行星之间的区别\(^\text{[216]}\)。
\item \textbf{Rosetta.} “罗塞塔号”（Rosetta）探测器对丘留莫夫－格拉西缅科彗星（Comet Churyumov–Gerasimenko）进行了轨道飞行。2014 年 11 月 12 日，其着陆器“菲莱号”（Philae）成功登陆彗星表面——这是历史上首次有航天器成功在彗星上着陆\(^\text{[217]}\)。
\end{itemize}
\subsection{分类}
天文学家可以根据控制彗星近日点与远日点的行星，以及其 Tisserand 轨道参数，对彗星进行分类[218]。此外，也存在其他更具定性或非正式的分类方式。
\subsubsection{大彗星}
\begin{figure}[ht]
\centering
\includegraphics[width=8cm]{./figures/5be89116e8296585.png}
\caption{1577 年大彗星的木刻版画} \label{fig_Comet_23}
\end{figure}
大约每十年会有一颗彗星变得足够明亮，使普通观察者也能注意到，这类彗星被称为“大彗星”[154]。预测一颗彗星是否会成为大彗星是出了名的困难，因为许多因素可能导致其亮度与预测值大幅偏离[219]。一般来说，如果一颗彗星拥有巨大且活跃的彗核、在近日点时非常接近太阳、且在最亮时不会因与太阳方向重合而从地球上难以观测，那么它就有成为大彗星的可能。然而，1973 年的科霍特克彗星（Comet Kohoutek）满足所有这些条件，本被期望成为壮观的彗星，却最终未能达到预期[220]。相比之下，三年后出现的 West 大彗星虽然最初预期不高，但却成为极为壮丽的大彗星[221]。

1577 年大彗星是大彗星的著名例子。它作为一颗非周期彗星，近距离掠过地球，被包括著名天文学家第谷·布拉赫（Tycho Brahe）与塔基丁（Taqi ad-Din）在内的众多观测者看到。对这颗彗星的观测促成了彗星科学的多项重要发现，尤其是对布拉赫而言。

20 世纪末出现了一个长期缺乏大彗星的空白期，随后两颗大彗星接连出现——1996 年的百武彗星（Comet Hyakutake），以及两年后达到最亮的海尔–波普彗星（Hale–Bopp），后者在 1995 年被发现。21 世纪的第一颗大彗星是 C/2006 P1（McNaught），其在 2007 年 1 月变得肉眼可见，是 40 多年来最明亮的彗星[222]。
\subsubsection{掠日彗星}  
掠日彗星指在近日点极其接近太阳的彗星，通常距离仅为几百万公里[223]。小型掠日彗星在如此接近太阳时可能会完全蒸发，而大型掠日彗星则可能经历多次近日点通过。然而，它们所承受的强大潮汐力常常会导致其发生裂解[224]。

在 SOHO（太阳和日光层观察卫星）的观测中，大约 90\% 的掠日彗星属于克罗伊茨族（Kreutz group），这一族群被认为源自一颗巨型彗星，它在第一次进入内太阳系时分裂成众多较小彗星[225]。其余掠日彗星中包含一些零散的成员，但人们在其中识别出另四个相关族群：Kracht、Kracht 2a、Marsden 和 Meyer 族群。其中 Marsden 与 Kracht 两族似乎均与 96P/Machholz 彗星有关，而后者又是象限仪（Quadrantids）与白羊座流星群（Arietids）这两条流星雨母体的来源[226]。
\subsubsection{异常彗星}
\begin{figure}[ht]
\centering
\includegraphics[width=8cm]{./figures/46cf07ad5d90f072.png}
\caption{展示太阳系天体类型的欧拉图} \label{fig_Comet_24}
\end{figure}
在已知的数千颗彗星中，有一些呈现出不寻常的性质。恩克彗星（Comet Encke，2P/Encke）的轨道从小行星带之外延伸到水星轨道以内，而施瓦斯曼–瓦赫曼 29 号彗星（Comet 29P/Schwassmann–Wachmann）目前则在木星和土星轨道之间以几乎圆形的轨道运行[227]。2060 Chiron 的不稳定轨道位于土星与天王星之间，起初被分类为小行星，直到后期观测到其具有微弱彗发[228]。同样地，舒梅克–列维 2 号彗星（Comet Shoemaker–Levy 2）最初被命名为小行星 1990 UL3[229]。
\subsubsection{最大彗星}  
已知最大的一颗周期彗星是 95P/Chiron，其直径约为 200 公里，每 50 年经过一次近日点，并在距离 8 AU、略低于土星轨道的位置通过。已知最大的一颗奥尔特云彗星被认为是伯纳迪内利–伯恩斯坦彗星（Comet Bernardinelli–Bernstein），直径约为 150 公里，将在 2031 年 1 月经过近日点，位置在距太阳 11 AU、略高于土星轨道之外。1729 年大彗星（Comet of 1729）的直径估计约为 100 公里，其近日点位于距太阳 4 AU、木星轨道以内。
\subsubsection{半人马天体}  
半人马天体通常同时表现出小行星与彗星的特征[230]。一些半人马天体被分类为彗星，如 60558 Echeclus 与 166P/NEAT。166P/NEAT 在发现时具有彗发，因此尽管其轨道特性不像典型彗星，仍被归类为彗星；而 60558 Echeclus 在被发现时没有彗发，但后来变得活跃[231]，因此被同时分类为彗星与小行星（174P/Echeclus）。曾有一个计划打算让卡西尼号（Cassini）飞往一颗半人马天体，但 NASA 最终决定将其销毁[232]。
\subsection{观测}
彗星可以通过使用广角望远镜的摄影方式发现，也可以通过双筒望远镜进行目视发现。然而，即使没有光学设备，业余天文学家仍可以通过下载由某些卫星天文台（如 SOHO）拍摄的累积图像，在网上发现掠日彗星[233]。SOHO 的第 2000 颗彗星由波兰业余天文学家 Michał Kusiak 于 2010 年 12 月 26 日发现[234]，而海尔–波普彗星的两位发现者也都使用了业余设备（尽管 Hale 本人不是业余天文学家）。
\subsubsection{遗失彗星}  
一些在早期几十年或几个世纪被发现的周期彗星如今已成为遗失彗星。由于其轨道从未被准确确定，以至无法预测其未来出现时间，或彗星已完全解体。然而，有时一颗“新”彗星被发现后，通过计算其轨道会显示它实际上是一颗古老的“遗失”彗星。例如，第 11P/Tempel–Swift–LINEAR 彗星于 1869 年被发现，但由于受到木星的扰动，自 1908 年后便无法观测。直至 2001 年，它才被 LINEAR 项目意外重新发现[235]。至少有 18 颗彗星属于这一类别[236]。
\subsection{在大众文化中} 
流行文化中对彗星的描绘深受西方长期传统影响，即将彗星视为厄运的预兆、世界变动的象征[237]。仅哈雷彗星就曾在每一次回归中引发大量耸动的出版物。尤其是一些名人的生与死恰逢彗星出现，例如作家马克·吐温（Mark Twain，曾正确预言自己将在 1910 年“随彗星而去”）[237]，以及作家尤多拉·韦尔蒂（Eudora Welty），Mary Chapin Carpenter 曾以歌曲《Halley Came to Jackson》献给她[237]。

在过去，明亮的彗星常使大众陷入恐慌与歇斯底里，被视为凶兆。较近时期，例如 1910 年哈雷彗星经过时，地球穿过其彗尾，而报纸上错误的报道引发了公众对彗尾中氰化物将毒害数百万人的恐惧[238]；1997 年海尔–波普彗星的出现则触发了“天堂之门”邪教的大规模集体自杀事件[239]。

在科幻作品中，彗星撞击常被描绘成可通过科技与英雄主义克服的威胁（如 1998 年的电影《天地大冲撞》（Deep Impact）与《世界末日》（Armageddon）），或是全球性灾难的起因（如 1979 年小说《路西法之锤》（Lucifer's Hammer）），或僵尸灾难的源头（如 1984 年电影《彗星之夜》（Night of the Comet））[237]。在儒勒·凡尔纳的《乘彗星旅行》（Off on a Comet）中，一群人被困在绕太阳运行的彗星上，而在阿瑟·克拉克的小说《2061：太空漫游三》（2061: Odyssey Three）中，一支大型载人太空探险队访问了哈雷彗星[240]。
\subsection{在文学中}
1825 年 7 月 15 日，庞斯（Pons）在佛罗伦萨首次记录的一颗长周期彗星，启发了 Lydia Sigourney 创作其幽默诗《1825 年的彗星》（The Comet of 1825），诗中描绘了各个天体就彗星的出现与目的展开争论。
\subsection{图集}
\begin{figure}[ht]
\centering
\includegraphics[width=14.25cm]{./figures/694e567e802d8ebd.png}
\caption{} \label{fig_Comet_25}
\end{figure}
\subsection{参见}
\begin{itemize}
\item The Big Splash  
\item 彗星年份（Comet vintages）  
\item 地球撞击坑列表（List of impact craters on Earth）  
\item 地球上的可能撞击构造列表（List of possible impact structures on Earth）  
\item 彗星列表（Lists of comets） 
\end{itemize} 
\subsection{参考文献}
\subsubsection{脚注}\\
a.“我并不认为彗星只是一团突然出现的火焰，而是自然界永恒作品的一部分。”（Sagan \& Druyan 1997, p.~26）\\
b.塞涅卡（Seneca）被引用说过：“为何……我们会惊讶于彗星——宇宙中如此罕见的景象——尚未被固定规律所掌握，其起始与终结也无人知晓？它们的回归间隔如此漫长……时机终将到来，经过极长时期的勤勉研究，人们会揭示那些如今仍被隐藏的事物。”[178]

Citations —— 中文翻译**

\begin{enumerate}
\item “问一位天文学家（Ask an Astronomer）”。Cool Cosmos。检索于 2023 年 3 月 11 日。
\item Randall, Lisa（2015）。《暗物质与恐龙：宇宙惊人的相互关联性》（Dark Matter and the Dinosaurs: The Astounding Interconnectedness of the Universe）。纽约：Ecco/HarperCollins Publishers。第 104–105 页。ISBN 978-0-06-232847-2。
\item “小行星与彗星的区别是什么（What is the difference between asteroids and comets）”。Rosetta 常见问题（European Space Agency）。检索于 2013 年 7 月 30 日。
\item “什么是小行星和彗星（What Are Asteroids And Comets）”。近地天体计划 FAQ（NASA）。原文于 2004 年 6 月 28 日存档。检索于 2013 年 7 月 30 日。
\item Ishii, H. A.; 等（2008）。《彗星 81P/Wild 2 的尘埃与来自彗星的行星际尘埃的比较》。Science，319（5862）：447–450。Bibcode:2008Sci...319..447I。doi:10.1126/science.1150683。PMID 18218892。S2CID 24339399。
\item “JPL 小天体数据库浏览器 C/2014 S3（PANSTARRS）”。
\item Stephens, Haynes; 等（2017 年 10 月）。《追逐 Manx 彗星：无尾的长周期彗星》（Chasing Manxes: Long-Period Comets Without Tails）。AAS/行星科学部会议摘要，49（49），420.02。Bibcode:2017DPS....4942002S。
\item “已发现的彗星（Comets Discovered）”。小行星中心。检索于 2021 年 4 月 27 日。
\item Erickson, Jon（2003）。《小行星、彗星与陨石：地球的宇宙入侵者》（Asteroids, Comets, and Meteorites: Cosmic Invaders of the Earth）。The Living Earth。纽约：Infobase。第 123 页。ISBN 978-0-8160-4873-1。
\item Couper, Heather; 等（2014）。《行星：太阳系权威指南》（The Planets: The Definitive Guide to Our Solar System）。伦敦：Dorling Kindersley。第 222 页。ISBN 978-1-4654-3573-6。

\item Licht, A.（1999）。《肉眼可见彗星的出现率：公元前 101 年至公元 1970 年》（The Rate of Naked-Eye Comets from 101 BC to 1970 AD）。Icarus，137（2）：355–356。Bibcode:1999Icar..137..355L。doi:10.1006/icar.1998.6048。

\item “着陆！Rosetta 的 Philae 探测器成功登陆彗星”。欧洲空间局。2014 年 11 月 12 日。检索于 2017 年 12 月 11 日。

\item Mardon, A. D. A. A. “第 65 届陨石学会年会（2002）5124”。[https://www.lpi.usra.edu/meetings/metsoc2002/pdf/5124.pdf](https://www.lpi.usra.edu/meetings/metsoc2002/pdf/5124.pdf)

\item Schove, D. J.（1975）。《公元 200–1882 年的彗星年表（数字形式）》(Comet chronology in numbers, AD 200–1882)。英国天文学会杂志（Journal of the British Astronomical Association），第 85 卷，第 401–407 页。

\item “comet”。《牛津英语词典》在线版（Oxford English Dictionary）。牛津大学出版社。（需要订阅或机构访问权限。）

\item Harper, Douglas. “Comet（n.）”。《在线词源词典》。检索于 2013 年 7 月 30 日。

\item 《大美百科全书》（The Encyclopedia Americana: A Library of Universal Knowledge）。第 26 卷。The Encyclopedia Americana Corp. 1920。第 162–163 页。

\item Greenberg, J. Mayo（1998）。《制造一个彗星核》（Making a comet nucleus）。Astronomy & Astrophysics，330：375。Bibcode:1998A&A...330..375G。

\item “太空中的脏雪球（Dirty Snowballs in Space）”。Starryskies。原文于 2013 年 1 月 29 日存档。检索于 2013 年 8 月 15 日。

\item “来自 ESA Rosetta 的证据表明：彗星与其说是‘脏雪球’，不如说更像‘冰质尘球’”。Times Higher Education。2005 年 10 月 21 日。

\item Clavin, Whitney（2015 年 2 月 10 日）。《为什么彗星像炸冰淇淋》（Why Comets Are Like Deep Fried Ice Cream）。NASA。检索于 2015 年 2 月 10 日。

\item Meech, M.（1997 年 3 月 24 日）。《1997 年彗星海尔–波普：我们能从明亮彗星中学到什么》。Planetary Science Research Discoveries。检索于 2013 年 4 月 30 日。

\item “Stardust 的结果表明彗星比想象中更复杂”。NASA。2006 年 12 月 14 日。检索于 2013 年 7 月 31 日。

\item Elsila, Jamie E.; 等（2009）。《在 Stardust 返回的样品中检测到彗星甘氨酸》（Cometary glycine detected in samples returned by Stardust）。Meteoritics & Planetary Science，44（9）：1323。Bibcode:2009M&PS...44.1323E。doi:10.1111/j.1945-5100.2009.tb01224.x。

\item Callahan, M. P.; 等（2011）。《碳质陨石含有多种地外核碱基》（Carbonaceous meteorites contain a wide range of extraterrestrial nucleobases）。PNAS，108（34）：13995–8。Bibcode:2011PNAS..10813995C。doi:10.1073/pnas.1106493108。PMC 3161613。PMID 21836052。

\item Steigerwald, John（2011 年 8 月 8 日）。《NASA 研究者：DNA 构件可以在太空中形成》。NASA。原文于 2020 年 4 月 26 日存档。检索于 2013 年 7 月 31 日。

\item Weaver, H. A.; 等（1997）。《海尔–波普彗星核的活动与大小》（The Activity and Size of the Nucleus of Comet Hale–Bopp）。Science，275（5308）：1900–1904。Bibcode:1997Sci...275.1900W。doi:10.1126/science.275.5308.1900。PMID 9072959。

\item Hanslmeier, Arnold（2008）。《宜居性与宇宙灾变》（Habitability and Cosmic Catastrophes）。Springer。第 91 页。ISBN 978-3-540-76945-3。

\item Fernández, Yanga R.（2000）。《海尔–波普彗星核：大小与活动》（The Nucleus of Comet Hale–Bopp）。Earth, Moon, and Planets，89（1）：3–25。Bibcode:2002EM&P...89....3F。

\item Jewitt, David（2003 年 4 月）。《彗星的核》（The Cometary Nucleus）。加州大学洛杉矶分校地球与空间科学系。检索于 2013 年 7 月 31 日。
\end{enumerate}










